% ---------------------------------------------------------------------
%  TÓM TẮT
% ---------------------------------------------------------------------
\begin{abstract}

Điểm danh thủ công ở giảng đường đại học vừa tốn thời gian của mỗi tiết
học, vừa không ngăn được tình trạng điểm danh hộ, và số liệu chuyên cần
thường chỉ tổng hợp được vào cuối kỳ, khi đã quá muộn để nhắc nhở sinh viên
sắp chạm ngưỡng cấm thi. Đồ án xây dựng FAMS \en{(Face Attendance Management
System)}, một hệ thống quản lý sinh viên tích hợp điểm danh bằng nhận diện
khuôn mặt, dùng được cho cả nhà trường lẫn sinh viên và phụ huynh.

Hệ thống nhận diện khuôn mặt theo hướng so khớp đặc trưng: bộ dò SCRFD tìm
khuôn mặt trong khung hình, mô hình ArcFace chuyển mỗi khuôn mặt thành vector
\FchieuVector{} chiều đã chuẩn hoá, và danh tính được xác định bằng khoảng
cách cosine tới các vector đã đăng ký, có chỉ mục HNSW của pgvector hỗ trợ
tìm kiếm. Cách này cho phép thêm một sinh viên mới chỉ bằng một ảnh, không
phải huấn luyện lại mô hình. Để chống gian lận, hệ thống kiểm tra ảnh in và
ảnh chụp lại màn hình bằng MiniFASNetV2, yêu cầu nhận diện đúng
\FsoKhungXacNhan{} khung hình mới ghi nhận, và đối chiếu toạ độ GPS khi sinh
viên tự điểm danh bằng ảnh tự chụp.

Sản phẩm gồm một dịch vụ FastAPI với \FsoEndpoint{} endpoint trên
\FsoBang{} bảng dữ liệu PostgreSQL, một ứng dụng Flutter dùng chung cho Web,
Android và iOS với \FsoManHinh{} màn hình, phân quyền cho \FsoVaiTro{} vai
trò từ quản trị viên, phòng ban, giảng viên đến sinh viên và phụ huynh. Bộ
kiểm thử tự động gồm \FsoTest{} ca chạy trên PostgreSQL thật. Kết quả cho
thấy một hệ thống điểm danh bằng khuôn mặt có thể chạy trên máy chủ không
GPU, đồng thời khẳng định rằng phần khó nhất của bài toán không nằm ở mô
hình nhận diện, mà ở các lớp phòng vệ quanh nó.

\vspace{0.8em}
\noindent\textbf{Từ khoá:} nhận diện khuôn mặt, ArcFace, chống giả mạo
khuôn mặt, tìm kiếm vector, điểm danh tự động.

\end{abstract}
